\documentclass{article}
\usepackage{neurips_2025}

\usepackage[utf8]{inputenc}
\usepackage[T1]{fontenc}
\usepackage{hyperref}
\usepackage{url}
\usepackage{booktabs}
\usepackage{amsfonts}
\usepackage{nicefrac}
\usepackage{microtype}
\usepackage{xcolor}
\usepackage{graphicx}

% ============================================================
%  TODO (Member 4): finalise title
% ============================================================
\title{Binary Medical Visual Question Answering via Multimodal Fusion of Image and Text Encoders}

\author{%
  % TODO (Member 1): fill in your name / email
  Member 1 Name \\
  Department of Computer Science \\
  University of Houston \\
  \texttt{member1@example.com} \\
  \And
  Phuong Ho \\
  Department of Computer Science \\
  University of Houston \\
  \texttt{duyphuong1799@gmail.com} \\
  \And
  % TODO (Member 3): fill in your name / email
  Member 3 Name \\
  Department of Computer Science \\
  University of Houston \\
  \texttt{member3@example.com} \\
  \And
  % TODO (Member 4): fill in your name / email
  Member 4 Name \\
  Department of Computer Science \\
  University of Houston \\
  \texttt{member4@example.com} \\
}

\begin{document}

\maketitle

% ============================================================
%  ABSTRACT — TODO (Member 4): write final version, max 300 words
%  Suggested flow: problem → dataset → approach → key results
% ============================================================
\begin{abstract}
Medical Visual Question Answering (Med-VQA) requires simultaneous understanding
of clinical images and natural-language questions, making it a challenging
multimodal classification task. In this work we address the binary (yes/no)
subset of Med-VQA using the \texttt{robailleo/medical-vision-llm-dataset}, which
contains mixed medical imaging modalities including X-ray, MRI, CT, and
ultrasound. We build a two-stage pipeline: a routing classifier that directs
questions to either a binary classification branch or an open-ended captioning
branch, followed by modality-specific models for each branch. For binary VQA we
systematically compare text-only baselines (TF-IDF with Logistic Regression and
SVM), an image-only CNN baseline (ResNet-18 with partial fine-tuning), and
multimodal fusion models combining ResNet-18 with either CLIP or BioBERT. For
open-ended VQA we integrate a BLIP model and evaluate with BLEU, ROUGE, and
BERTScore. Our best binary classifier---ResNet-18 fused with a fine-tuned
BioBERT---achieves \textbf{79.5\% accuracy} and \textbf{0.83 AUC-ROC},
outperforming single-modality baselines by over 7 percentage points.
% TODO (Member 4): update numbers and add captioning results once confirmed
\end{abstract}


% ============================================================
%  1. INTRODUCTION — TODO (Member 4): write 3--4 paragraphs
%  Para 1: motivate Med-VQA — clinical decision support
%  Para 2: challenge — multimodal, small labeled data, mixed modalities
%  Para 3: our pipeline overview and key contributions
%  Para 4: roadmap of the paper
% ============================================================
\section{Introduction}

Medical Visual Question Answering (VQA) has attracted significant research
interest as a potential tool for automated clinical decision support. A Med-VQA
system must jointly reason over a clinical image and a natural-language question
to produce a correct answer. Questions may be binary (``Is there a fracture?'')
or open-ended (``What abnormality is visible?''), requiring different modelling
strategies.

% TODO (Member 4): write the remaining introduction paragraphs here.


% ============================================================
%  2. BACKGROUND — TODO (Member 4): related work, 1--2 pages
%  Cover: (1) VQA overview, (2) Med-VQA prior work, (3) transfer learning in
%         medical imaging, (4) biomedical language models, (5) image captioning
%  Cite at least 6--8 relevant papers
% ============================================================
\section{Background}

\subsection{Visual Question Answering}
% TODO (Member 4): overview of VQA task and key datasets/models (VQA v2, etc.)

\subsection{Medical VQA}
% TODO (Member 4): VQA-Med, PathVQA, SLAKE — datasets and prior systems

\subsection{Transfer Learning for Medical Imaging}
% TODO (Member 4): ImageNet pretraining, CheXNet, domain generalisation

\subsection{Biomedical Language Models}
% TODO (Member 4): BERT, BioBERT, PubMedBERT — pretraining and downstream tasks

\subsection{Vision-Language Models and Image Captioning}
% TODO (Member 4): CLIP, BLIP, GIT — and their application to medical imaging


% ============================================================
%  3. DATA — Owner: Member 1
% ============================================================
\section{Data}

% TODO (Member 1): write this entire section based on your EDA work

\subsection{Dataset Description}
% TODO (Member 1): describe robailleo/medical-vision-llm-dataset
%   - source, modalities (X-ray, MRI, CT, ultrasound), total size
%   - answer type breakdown (binary vs open-ended, class distribution)
%   - image size distribution

We use the \texttt{robailleo/medical-vision-llm-dataset} available on
HuggingFace~\cite{TODO-dataset}.
% TODO (Member 1): fill in dataset statistics here

\subsection{Preprocessing Pipeline}
% TODO (Member 1): describe your full preprocessing pipeline
%   - resizing to 224x224
%   - grayscale to 3-channel RGB conversion
%   - ImageNet normalisation (mean/std)
%   - tokenisation strategy for text

Images are resized to $224 \times 224$ pixels and normalized using ImageNet
statistics (mean $[0.485, 0.456, 0.406]$, std $[0.229, 0.224, 0.225]$).
Grayscale images are converted to three-channel RGB. Questions are tokenized
with a maximum sequence length of 128 tokens.

\subsection{Train/Validation Split}
% TODO (Member 1): describe split sizes, stratification strategy, class balance
%   - report exact numbers after filtering to binary subset


% ============================================================
%  4. METHODS
% ============================================================
\section{Methods}

Our system consists of three components: (1) a routing classifier that
separates binary from open-ended questions, (2) a binary VQA pipeline, and
(3) an open-ended captioning pipeline. Figure~\ref{fig:pipeline} illustrates
the overall architecture.

% TODO (Member 4): add a pipeline diagram figure here
% \begin{figure}[h]
%   \centering
%   \includegraphics[width=0.9\linewidth]{figures/pipeline.png}
%   \caption{End-to-end pipeline: routing classifier directs inputs to the
%            binary classification or open-ended captioning branch.}
%   \label{fig:pipeline}
% \end{figure}

\subsection{Routing Classifier (Member 2)}
To separate binary yes/no questions from open-ended questions, we train a
lightweight routing classifier on the question text. We extract TF-IDF features
from the question and train a Support Vector Machine (SVM) classifier to predict
whether the expected answer is binary or open-ended. This allows each downstream
model to specialise on its task.
% TODO (Member 2 / Phuong): add any additional details about the routing classifier
%   - vocabulary size, TF-IDF settings, SVM kernel, accuracy on routing task

\subsection{Text-Only Baselines (Member 2)}
For binary VQA, we establish text-only baselines using TF-IDF feature
extraction combined with Logistic Regression and Support Vector Machine
classifiers trained only on question text, ignoring the image entirely. These
baselines reveal how much information is carried by question phrasing alone.
% TODO (Member 2 / Phuong): describe feature extraction details and hyperparameters

\subsection{Image-Only Baseline (Member 2)}
We fine-tune a ResNet-18~\cite{he2016resnet} pretrained on ImageNet. All layers
except the final residual block (\texttt{layer4}) are frozen; only
\texttt{layer4} and a new binary classification head are trained. This strategy
reuses low-level ImageNet features while adapting high-level representations to
the medical imaging domain.

\subsection{Multimodal Fusion Models (Member 2)}

\subsubsection{Architecture}
Our fusion model concatenates image and text feature vectors before a learned
classification head:
\begin{itemize}
    \item \textbf{Image encoder}: ResNet-18 with \texttt{layer4} unfrozen,
          producing a 512-dim feature vector after global average pooling.
    \item \textbf{Text encoder}: Either a frozen CLIP text encoder (512-dim)
          or BioBERT (\texttt{dmis-lab/biobert-base-cased-v1.2}~\cite{lee2020biobert},
          768-dim) with mean pooling over the last hidden state.
    \item \textbf{Fusion head}: LayerNorm on each stream $\rightarrow$
          concatenation (1024-dim or 1280-dim) $\rightarrow$
          Linear(256)--ReLU--Dropout(0.5)--Linear(1).
\end{itemize}

\subsubsection{Why BioBERT over CLIP}
CLIP~\cite{radford2021clip} is pretrained on internet image--caption pairs and
uses general-purpose language. BioBERT~\cite{lee2020biobert} is pretrained on
PubMed abstracts and PubMed Central full texts, making it significantly better
suited to understanding medical terminology and clinical question phrasing.

\subsubsection{Encoder Fine-Tuning}
Following the professor's suggestion that domain-adapted fine-tuning may
further improve performance, we additionally fine-tune the top two transformer
layers of BioBERT (layers 10 and 11) at a very low learning rate
($5 \times 10^{-6}$) to adapt the text representations to our specific dataset.

\subsubsection{Training Details}
Models are trained with \texttt{BCEWithLogitsLoss} and AdamW with separate
learning rate groups: $10^{-4}$ for the fusion head and normalisation layers,
$10^{-5}$ for ResNet \texttt{layer4}, and $5 \times 10^{-6}$ for fine-tuned
BioBERT layers. Batch size 16, weight decay $10^{-4}$, up to 30 epochs,
early stopping with patience 7, and \texttt{ReduceLROnPlateau} scheduling
(factor 0.5, patience 3).

\subsection{Vision-Language Model: BLIP (Member 3)}
% TODO (Member 3): write this entire subsection
%   - BLIP model description and why you chose it
%   - how you froze the vision encoder and fine-tuned the language head
%   - input format (image + prompt)
%   - evaluation metrics: BLEU, ROUGE, BERTScore
%   - any prompt engineering experiments

For the open-ended captioning branch we integrate
BLIP~\cite{TODO-blip} % TODO: add BLIP citation
as a vision-language model capable of generating free-form answers given an
image and question prompt.

% TODO (Member 3): fill in full details here


% ============================================================
%  5. EXPERIMENTS — Owner: Member 4 (tables, figures, analysis)
%     Each member contributes their own results; Member 4 assembles
% ============================================================
\section{Experiments}

\subsection{Evaluation Metrics}
For binary VQA we report accuracy, macro F1 score, and AUC-ROC on the
validation split. For open-ended captioning we report BLEU-4, ROUGE-L, and
BERTScore~\cite{TODO-bertscore}.
% TODO (Member 4): add citations for metrics

\subsection{Binary VQA Results}

Table~\ref{tab:binary-results} summarises the performance of all binary VQA
models.

\begin{table}[h]
  \caption{Binary VQA model comparison on the validation set.}
  \label{tab:binary-results}
  \centering
  \begin{tabular}{lccc}
    \toprule
    Model & Accuracy & F1 & AUC-ROC \\
    \midrule
    TF-IDF + Logistic Regression             & 0.6684 & 0.6441 & 0.7073 \\
    TF-IDF + SVM                             & 0.6895 & 0.6704 & 0.7517 \\
    ResNet-18 (layer4 unfrozen)              & 0.7263 & 0.7500 & 0.7740 \\
    Fusion: ResNet-18 + CLIP                 & 0.7579 & 0.7723 & 0.8194 \\
    Fusion: ResNet-18 + BioBERT              & 0.7789 & 0.7835 & 0.8204 \\
    Fusion: ResNet-18 + BioBERT (fine-tuned) & \textbf{0.7947} & \textbf{0.8000} & \textbf{0.8281} \\
    \bottomrule
  \end{tabular}
\end{table}

\subsection{Open-Ended Captioning Results (Member 3)}
% TODO (Member 3): provide your BLEU/ROUGE/BERTScore numbers
% TODO (Member 4): format into a table

\begin{table}[h]
  \caption{Open-ended captioning results on the validation set.}
  \label{tab:captioning-results}
  \centering
  \begin{tabular}{lccc}
    \toprule
    Model & BLEU-4 & ROUGE-L & BERTScore \\
    \midrule
    % TODO (Member 3): fill in rows
    BLIP (frozen encoder)       & TODO & TODO & TODO \\
    BLIP (fine-tuned lang head) & TODO & TODO & TODO \\
    \bottomrule
  \end{tabular}
\end{table}

\subsection{Ablation Study (Member 2)}
% TODO (Member 2 / Phuong): discuss the effect of each design choice:
%   - text encoder: CLIP vs BioBERT (+2.1 pp accuracy)
%   - BioBERT fine-tuning: frozen vs layers 10+11 unfrozen (+1.6 pp)
%   - LayerNorm: with vs without (if you have data)
%   - CheXNet as image encoder: why it failed (domain mismatch)

\subsection{End-to-End Pipeline Evaluation (Member 4)}
% TODO (Member 4): describe how the routing classifier connects to the two
%   branches; report routing accuracy and overall end-to-end metrics

\subsection{Qualitative Analysis (Member 4)}
% TODO (Member 4): include saliency / attention map visualisations
%   - show examples where the model is correct vs incorrect
%   - ROC curve figure

% \begin{figure}[h]
%   \centering
%   \includegraphics[width=0.8\linewidth]{figures/roc_curve.png}
%   \caption{ROC curves for binary VQA models.}
%   \label{fig:roc}
% \end{figure}

\subsection{Compute Resources}
% TODO (Member 4): GPU type (Colab T4/A100), approx. training time per model


% ============================================================
%  6. CONCLUSION — TODO (Member 4)
% ============================================================
\section{Conclusion}

We presented a two-stage pipeline for Medical VQA that routes questions to
either a binary classification branch or an open-ended captioning branch. For
binary questions, combining a ResNet-18 image encoder with a fine-tuned BioBERT
text encoder achieves 79.5\% accuracy and 0.83 AUC-ROC, outperforming all
single-modality baselines. The consistent advantage of BioBERT over CLIP
confirms the value of domain-adapted language models for medical text.
% TODO (Member 4): add captioning conclusions once Member 3's results are in

\paragraph{Future work.}
% TODO (Member 4): suggest 2--3 directions, e.g.:
%   - larger and more diverse medical VQA datasets
%   - cross-attention fusion instead of simple concatenation
%   - end-to-end trainable routing via multi-task learning


% ============================================================
\begin{ack}
This project was completed as part of COSC 4368 at the University of Houston.
% TODO (Member 4): add instructor/TA name if appropriate
\end{ack}


% ============================================================
\section*{References}
{\small
% TODO (Member 4): replace TODOs with full citations; use consistent style

[1] He, K., Zhang, X., Ren, S., \& Sun, J. (2016). Deep residual learning for image recognition. \textit{CVPR}.

[2] Lee, J., Yoon, W., Kim, S., et al. (2020). BioBERT: a pre-trained biomedical language representation model for biomedical text mining. \textit{Bioinformatics}, 36(4), 1234--1240.

[3] Radford, A., Kim, J.W., Hallacy, C., et al. (2021). Learning transferable visual models from natural language supervision. \textit{ICML}.

[4] Devlin, J., Chang, M., Lee, K., \& Toutanova, K. (2019). BERT: Pre-training of deep bidirectional transformers for language understanding. \textit{NAACL}.

% TODO (Member 4): add citations for:
%   - BLIP (Li et al. 2022)
%   - BERTScore (Zhang et al. 2020)
%   - VQA-Med dataset
%   - robailleo/medical-vision-llm-dataset
%   - any other related work cited in Background
}


% ============================================================
\appendix
\section{Supplementary Material}
% Not counted toward the 4--6 page limit

\subsection{Source Code}
% TODO (Member 4): update with actual repo URL
All notebooks and scripts are available at:
\url{https://github.com/TODO/Medical_VQA_Classification}

\subsection{Model Architecture Diagram}
% TODO (Member 4): add fusion architecture figure if not included in main paper

\subsection{Training Curves}
% TODO (Member 4): include loss/accuracy curves for key models


% ============================================================
\newpage
\section*{NeurIPS Paper Checklist}

\begin{enumerate}

\item {\bf Claims}
    \item[] Question: Do the main claims made in the abstract and introduction accurately reflect the paper's contributions and scope?
    \item[] Answer: \answerYes{}
    \item[] Justification: All reported numbers are taken directly from experimental results in Tables~\ref{tab:binary-results} and~\ref{tab:captioning-results}.

\item {\bf Limitations}
    \item[] Question: Does the paper discuss the limitations of the work performed by the authors?
    \item[] Answer: \answerYes{}
    \item[] Justification: We discuss the small dataset size as the primary bottleneck and the overfitting behaviour in Section~5.

\item {\bf Theory assumptions and proofs}
    \item[] Question: For each theoretical result, does the paper provide the full set of assumptions and a complete proof?
    \item[] Answer: \answerNA{}
    \item[] Justification: This paper presents empirical results only; no theoretical proofs are included.

\item {\bf Experimental result reproducibility}
    \item[] Question: Does the paper fully disclose all information needed to reproduce the main results?
    \item[] Answer: \answerYes{}
    \item[] Justification: Architecture, optimizer, hyperparameters, and random seed are reported in Section~4. Code is provided in the supplementary material.

\item {\bf Open access to data and code}
    \item[] Question: Does the paper provide open access to data and code?
    \item[] Answer: \answerYes{}
    \item[] Justification: The dataset is publicly available on HuggingFace. Code is linked in the appendix.

\item {\bf Experimental setting/details}
    \item[] Question: Does the paper specify all training and test details necessary to understand the results?
    \item[] Answer: \answerYes{}
    \item[] Justification: Optimizer, learning rates, batch size, epochs, and early stopping criteria are all reported in Section~4.4.4.

\item {\bf Experiment statistical significance}
    \item[] Question: Does the paper report error bars or statistical significance tests?
    \item[] Answer: \answerNo{}
    \item[] Justification: Due to computational constraints, each model was run once. Error bars were not computed.

\item {\bf Experiments compute resources}
    \item[] Question: Does the paper provide information on compute resources?
    \item[] Answer: \answerYes{}
    \item[] Justification: GPU type and approximate training time are reported in Section~5.5.

\item {\bf Code of ethics}
    \item[] Question: Does the research conform with the NeurIPS Code of Ethics?
    \item[] Answer: \answerYes{}
    \item[] Justification: This work uses a publicly available dataset and does not involve human subjects or sensitive data collection.

\item {\bf Broader impacts}
    \item[] Question: Does the paper discuss societal impacts?
    \item[] Answer: \answerYes{}
    \item[] Justification: % TODO (Member 4): add a broader impacts sentence in the Conclusion

\item {\bf Safeguards}
    \item[] Question: Does the paper describe safeguards for responsible release?
    \item[] Answer: \answerNA{}
    \item[] Justification: We release no new models or datasets; all assets used are existing public resources.

\item {\bf Licenses for existing assets}
    \item[] Question: Are existing assets properly credited?
    \item[] Answer: \answerYes{}
    \item[] Justification: The dataset and all pretrained models (ResNet-18, BioBERT, CLIP, BLIP) are cited in the references section.

\item {\bf New assets}
    \item[] Question: Are new assets well documented?
    \item[] Answer: \answerNA{}
    \item[] Justification: This paper introduces no new datasets or models.

\item {\bf Crowdsourcing and research with human subjects}
    \item[] Question: Does the paper involve crowdsourcing or human subjects research?
    \item[] Answer: \answerNA{}
    \item[] Justification: No human subjects were involved in this research.

\item {\bf Institutional review board (IRB) approvals}
    \item[] Question: Does the paper describe IRB approvals?
    \item[] Answer: \answerNA{}
    \item[] Justification: No human subjects research was conducted.

\item {\bf Declaration of LLM usage}
    \item[] Question: Does the paper describe usage of LLMs as a core method?
    \item[] Answer: \answerNA{}
    \item[] Justification: LLMs were not used as a core methodological component. BioBERT and BLIP are used as pretrained feature extractors/generators, not as novel LLM contributions.

\end{enumerate}

\end{document}
